\documentclass[11pt,a4paper]{article}

\usepackage{fontspec}
\usepackage{polyglossia}
\setmainlanguage{russian}
\setmainfont{Liberation Serif}
\setsansfont{Liberation Sans}
\setmonofont{DejaVu Sans Mono}

\usepackage{amsmath,amssymb}
\usepackage{booktabs}
\usepackage{geometry}
\geometry{margin=1.8cm}

\setlength{\parskip}{0.45em}
\setlength{\parindent}{0pt}

\begin{document}

\begin{center}
{\Large Ручной расчёт одной строки}\\[0.4em]
{\large Лабораторная работа №4, вариант 2}\\[0.2em]
Квадратуры Гаусса с тремя и четырьмя узлами
\end{center}

Сверяем программу \texttt{3\_gauss.py} с прямым счётом.
Берём ту же строку, что и в расчёте прямоугольников: \(j = 7\), \(t = 0{,}35\).
Тогда два метода можно положить рядом.

\section*{Что считаем}

Параметры варианта~2 те же:
\[
a = 0,\quad b = 2,\quad
t_{7} = 0{,}35,
\]
\[
f(x, t) = \cos\left(\frac{t}{1 + x^{2}} + 0{,}001\, x\right),
\qquad
F(0{,}35) = \int_{0}^{2} f(x,\, 0{,}35)\, dx.
\]
Буква \(t\) внутри интеграла постоянна. При \(t = 0{,}35\) функция одной переменной \(x\)
обозначена \(g(x) = f(x, 0{,}35)\). Дальше речь об интеграле от \(g\).

Число узлов фиксировано: отдельно 3 и отдельно 4.
Ниже сначала записаны определения, потом вывод узлов для \(n = 3\),
и только после этого численный счёт этой строки.

\section*{Квадратурная формула, узлы и веса}

Квадратурная формула заменяет интеграл конечной суммой:
\[
\int_{a}^{b} g(x)\, dx \approx \sum_{i=1}^{n} c_{i}\, g(x_{i}).
\]
Точки \(x_{i}\) называются \emph{узлами}: это те места, где вычисляют функцию.
Числа \(c_{i}\) называются \emph{весами}, или коэффициентами квадратуры.
Вес --- множитель при значении \(g(x_{i})\). Это не длина шага сетки.
Шаг появляется только у формул с равномерными узлами, например у прямоугольников.
У Гаусса узлы расставлены неравномерно, и у каждого узла свой вес.

Если узлы уже выбраны, веса можно подобрать так, чтобы формула была точна
для многочленов степени не выше \(n - 1\). Так устроены формулы Ньютона--Котеса:
прямоугольники, трапеции, Симпсон. Под интегралом функцию заменяют
её интерполяционным многочленом Лагранжа по этим узлам и интегрируют многочлен точно.
Свободны только \(n\) весов, поэтому хватает точности на степени \(0, 1, \ldots, n - 1\).

\section*{Алгебраическая степень точности}

\emph{Алгебраическая степень точности} формулы --- наибольшее целое \(m\),
при котором равенство
\[
\int_{a}^{b} x^{k}\, dx = \sum_{i=1}^{n} c_{i}\, x_{i}^{k}
\]
верно для всех \(k = 0, 1, \ldots, m\).
Если оно верно для этих степеней, оно верно для любого многочлена степени не выше \(m\):
многочлен есть сумма степеней, а интеграл и сумма линейны.

У формулы Ньютона--Котеса с \(n\) узлами степень точности равна \(n - 1\).
У Гаусса свободны и \(n\) узлов, и \(n\) весов, всего \(2n\) чисел.
Их выбирают так, чтобы выполнить \(2n\) условий: точность на степенях
\(0, 1, \ldots, 2n - 1\). Поэтому степень точности квадратуры Гаусса с \(n\) узлами равна \(2n - 1\).
Три узла точны для многочленов до степени 5, четыре узла --- до степени 7.

Косинус многочленом не является, так что для нашей \(g\) равенство приближённое.
Степень точности говорит, какие многочлены формула берёт без ошибки,
и за счёт этого оценивается порядок остатка для гладкой функции.

\section*{Ортогональные многочлены}

Два многочлена \(p\) и \(q\) называют \emph{ортогональными} на отрезке \([a, b]\) с весом \(w(x) > 0\), если
\[
\int_{a}^{b} p(x)\, q(x)\, w(x)\, dx = 0.
\]
Вес \(w\) --- функция, с которой их перемножают под интегралом.
В этой лабораторной на отрезке \([-1, 1]\) вес равен единице: \(w(x) = 1\).
Это другой «вес», не путать с коэффициентами \(c_{i}\) квадратуры.

\emph{Семейство ортогональных многочленов} --- последовательность \(P_{0}, P_{1}, P_{2}, \ldots\),
в которой степень \(P_{n}\) равна \(n\) и
\[
\int_{a}^{b} P_{m}(x)\, P_{n}(x)\, w(x)\, dx = 0
\quad\text{при } m \neq n.
\]
Отсюда следует более сильное свойство: \(P_{n}\) ортогонален \emph{любому} многочлену
степени меньше \(n\). Такой многочлен раскладывается по \(P_{0}, \ldots, P_{n-1}\),
а \(P_{n}\) ортогонален каждому из них.

\section*{Многочлены Лежандра}

\emph{Многочлены Лежандра} --- ортогональные многочлены на \([-1, 1]\) с весом \(w(x) = 1\).
Их задаёт формула Родрига:
\[
P_{n}(x) = \frac{1}{2^{n} n!}\, \frac{d^{n}}{dx^{n}} (x^{2} - 1)^{n}.
\]
Ниже выписаны \(P_{0}, \ldots, P_{4}\) в стандартной нормировке \(P_{n}(1) = 1\):
\[
\begin{aligned}
P_{0}(x) &= 1,\\
P_{1}(x) &= x,\\
P_{2}(x) &= \frac{3x^{2} - 1}{2},\\
P_{3}(x) &= \frac{5x^{3} - 3x}{2},\\
P_{4}(x) &= \frac{35x^{4} - 30x^{2} + 3}{8}.
\end{aligned}
\]
Квадрат нормы известен заранее:
\[
\int_{-1}^{1} P_{n}^{2}(x)\, dx = \frac{2}{2n + 1}.
\]

\section*{Почему узлы Гаусса --- корни Лежандра}

На \([-1, 1]\) квадратура Гаусса с \(n\) узлами имеет вид
\[
\int_{-1}^{1} g(x)\, dx \approx \sum_{i=1}^{n} c_{i}\, g(x_{i}).
\]
Узлы \(x_{i}\) берут корнями \(P_{n}\). Причина такая.

Пусть \(g\) --- многочлен степени не выше \(2n - 1\).
Поделим его на \(P_{n}\):
\[
g(x) = q(x)\, P_{n}(x) + r(x),
\]
где степень \(q\) не выше \(n - 1\), а степень остатка \(r\) тоже не выше \(n - 1\).
Интеграл произведения \(q P_{n}\) равен нулю: \(P_{n}\) ортогонален всякому многочлену
степени меньше \(n\), а \(q\) как раз такой. Значит
\[
\int_{-1}^{1} g(x)\, dx = \int_{-1}^{1} r(x)\, dx.
\]
Остаток \(r\) имеет степень не выше \(n - 1\). Любая квадратура с этими \(n\) узлами,
точная на многочленах степени не выше \(n - 1\), интегрирует \(r\) без ошибки.
В узлах \(P_{n}(x_{i}) = 0\), поэтому \(g(x_{i}) = r(x_{i})\).
Сумма по узлам для \(g\) совпадает с суммой для \(r\) и равна интегралу \(g\).

Так формула с \(n\) узлами становится точной на всех многочленах степени не выше \(2n - 1\).
Выше этой степени она уже не обязана быть точной: для этого не хватает свободных параметров.

\section*{Три узла: система из методички и многочлен \(P_{3}\)}

Для \(n = 3\) нужно шесть чисел: \(c_{1}, c_{2}, c_{3}, x_{1}, x_{2}, x_{3}\).
Их находят из точности на степенях \(k = 0, 1, \ldots, 5\):
\[
\int_{-1}^{1} x^{k}\, dx = \sum_{i=1}^{3} c_{i}\, x_{i}^{k}.
\]
Левая часть равна \(2/(k+1)\) при чётном \(k\) и равна нулю при нечётном \(k\).

Узлы Лежандра симметричны относительно нуля, и веса в симметричных узлах равны.
Для трёх узлов это значит
\[
x_{1} = -x_{3},\qquad x_{2} = 0,\qquad c_{1} = c_{3}.
\]
Условия при нечётных \(k = 1, 3, 5\) тогда выполняются сами.
Остаются три уравнения, они же выписаны в методичке:
\[
\begin{aligned}
k = 0:&\quad 2 = 2c_{1} + c_{2},\\
k = 2:&\quad \frac{2}{3} = 2c_{1} x_{1}^{2},\\
k = 4:&\quad \frac{2}{5} = 2c_{1} x_{1}^{4}.
\end{aligned}
\]
Делим третье на второе: \(x_{1}^{2} = 3/5\), так что
\[
x_{1} = -\sqrt{\frac{3}{5}},\qquad
x_{3} = \sqrt{\frac{3}{5}},\qquad
x_{2} = 0.
\]
Из второго уравнения
\[
\frac{2}{3} = 2c_{1}\cdot\frac{3}{5}
\quad\Rightarrow\quad
c_{1} = \frac{5}{9}.
\]
Из первого
\[
c_{2} = 2 - 2\cdot\frac{5}{9} = \frac{8}{9}.
\]
Итого
\[
c_{1} = c_{3} = \frac{5}{9},\qquad c_{2} = \frac{8}{9}.
\]

Те же узлы получаются как корни \(P_{3}\). Решаем \(P_{3}(x) = 0\):
\[
\frac{5x^{3} - 3x}{2} = 0
\quad\Rightarrow\quad
x(5x^{2} - 3) = 0
\quad\Rightarrow\quad
x = 0,\quad x = \pm\sqrt{\frac{3}{5}}.
\]
Система моментов и корни Лежандра дают один и тот же набор.

Для счёта на калькуляторе:
\[
\sqrt{\frac{3}{5}} = \sqrt{0{,}6} = 0{,}7745966692,
\quad
\frac{5}{9} = 0{,}5555555556,
\quad
\frac{8}{9} = 0{,}8888888889.
\]

\section*{Четыре узла}

При \(n = 4\) система уже на восемь неизвестных. В задании узлы и веса берут из таблицы.
Эти узлы --- корни \(P_{4}\), округлённые до шести знаков:
\[
\begin{aligned}
x_{1} &= -0{,}861136, & x_{2} &= -0{,}339981,\\
x_{3} &= 0{,}339981, & x_{4} &= 0{,}861136,\\
c_{1} &= c_{4} = 0{,}347855, & c_{2} &= c_{3} = 0{,}652145.
\end{aligned}
\]
Симметрия та же: узлы идут парами \(\pm x\), веса в паре равны.
Степень точности равна \(2\cdot 4 - 1 = 7\).

\section*{Остаток}

Если \(g\) имеет непрерывную производную порядка \(2n\), остаток формулы Гаусса на \([a, b]\) равен
\[
R_{n}(g)
= \frac{(b - a)^{2n + 1}\,(n!)^{4}}{(2n + 1)\,[(2n)!]^{3}}\, g^{(2n)}(\xi)
\]
при некотором \(\xi\) внутри \([a, b]\).
Для \(n = 3\) в остаток входит шестая производная, для \(n = 4\) --- восьмая.
В этой строке производные косинуса не считаем: практическая проверка --- совпадение
трёх узлов, четырёх узлов и прямоугольников в пределах \(\varepsilon = 0{,}001\).

\section*{Перенос на отрезок \([0, 2]\)}

В методичке для произвольного \([a, b]\):
\[
u_{i} = \frac{b - a}{2}\, x_{i} + \frac{b + a}{2},
\qquad
d_{i} = \frac{b - a}{2}\, c_{i},
\]
\[
\int_{a}^{b} g(u)\, du \approx \sum_{i=1}^{n} d_{i}\, g(u_{i}).
\]

У нас \(a = 0\), \(b = 2\), поэтому обе дроби равны единице:
\[
\frac{b - a}{2} = 1,
\qquad
\frac{b + a}{2} = 1.
\]
Формулы упрощаются:
\[
u_{i} = x_{i} + 1,
\qquad
d_{i} = c_{i}.
\]
Вес на \([0, 2]\) совпадает с табличным весом, узел просто сдвигается на 1 вправо.
Дальше в каждой точке считаем
\[
f(u) = \cos\left(\frac{0{,}35}{1 + u^{2}} + 0{,}001\, u\right)
\]
и складываем произведения \(d_{i} f(u_{i})\).

\section*{Три узла}

\[
\begin{aligned}
u_{1} &= 1 - 0{,}7745966692 = 0{,}2254033308, & d_{1} &= 0{,}5555555556,\\
u_{2} &= 1, & d_{2} &= 0{,}8888888889,\\
u_{3} &= 1 + 0{,}7745966692 = 1{,}7745966692, & d_{3} &= 0{,}5555555556.
\end{aligned}
\]

{\small
\begin{center}
\begin{tabular}{@{}r r r r r r@{}}
\toprule
\(u\) & \(1+u^{2}\) & аргумент & \(f(u)\) & \(d\) & \(d\,f(u)\) \\
\midrule
0,2254033308 & 1,0508066615 & 0,3333028502 & 0,9449669198 & 0,5555555556 & 0,5249816221 \\
1 & 2 & 0,176 & 0,9845519384 & 0,8888888889 & 0,8751572786 \\
1,7745966692 & 4,1491933385 & 0,0861283424 & 0,9962932466 & 0,5555555556 & 0,5534962481 \\
\bottomrule
\end{tabular}
\end{center}
}

Проверка средней строки, она считается без корня. \(u = 1\):
\[
1 + u^{2} = 2,
\qquad
\frac{0{,}35}{2} = 0{,}175,
\qquad
0{,}001\cdot 1 = 0{,}001,
\]
\[
0{,}175 + 0{,}001 = 0{,}176,
\qquad
\cos 0{,}176 = 0{,}9845519384,
\]
\[
\frac{8}{9}\cdot 0{,}9845519384 = 0{,}8751572786.
\]

Проверка первой строки, \(u = 0{,}2254033308\):
\[
u^{2} = 0{,}0508066615,
\quad
1 + u^{2} = 1{,}0508066615,
\quad
\frac{0{,}35}{1{,}0508066615} = 0{,}3330774469,
\]
\[
0{,}3330774469 + 0{,}0002254033 = 0{,}3333028502,
\]
\[
\cos 0{,}3333028502 = 0{,}9449669198,
\]
\[
0{,}5555555556 \cdot 0{,}9449669198 = 0{,}5249816221.
\]

Сумма трёх вкладов:
\[
0{,}5249816221 + 0{,}8751572786 + 0{,}5534962481 = 1{,}9536351488.
\]
Это и есть \(F\) по трём узлам.

\section*{Четыре узла}

Узлы сдвигаем на 1, веса оставляем табличными:
\[
\begin{aligned}
u_{1} &= 0{,}138864, & d_{1} &= 0{,}347855,\\
u_{2} &= 0{,}660019, & d_{2} &= 0{,}652145,\\
u_{3} &= 1{,}339981, & d_{3} &= 0{,}652145,\\
u_{4} &= 1{,}861136, & d_{4} &= 0{,}347855.
\end{aligned}
\]

{\small
\begin{center}
\begin{tabular}{@{}r r r r r r@{}}
\toprule
\(u\) & \(1+u^{2}\) & аргумент & \(f(u)\) & \(d\) & \(d\,f(u)\) \\
\midrule
0,138864 & 1,0192832105 & 0,3435174230 & 0,9415758208 & 0,347855 & 0,3275318571 \\
0,660019 & 1,4356250804 & 0,2444562613 & 0,9702690688 & 0,652145 & 0,6327561219 \\
1,339981 & 2,7955490804 & 0,1265389991 & 0,9920046180 & 0,652145 & 0,6469308516 \\
1,861136 & 4,4638272105 & 0,0802691889 & 0,9967801580 & 0,347855 & 0,3467349619 \\
\bottomrule
\end{tabular}
\end{center}
}

Проверка первой строки, \(u = 0{,}138864\):
\[
u^{2} = 0{,}0192832105,
\quad
1 + u^{2} = 1{,}0192832105,
\quad
\frac{0{,}35}{1{,}0192832105} = 0{,}3433785590,
\]
\[
0{,}3433785590 + 0{,}000138864 = 0{,}3435174230,
\]
\[
\cos 0{,}3435174230 = 0{,}9415758208,
\]
\[
0{,}347855 \cdot 0{,}9415758208 = 0{,}3275318571.
\]

Сумма четырёх вкладов:
\[
0{,}3275318571 + 0{,}6327561219 + 0{,}6469308516 + 0{,}3467349619 = 1{,}9539537925.
\]

\section*{Ответ и сверка}

\[
F_{3}(0{,}35) = 1{,}9536351488,
\qquad
F_{4}(0{,}35) = 1{,}9539537925.
\]

Строка программы \texttt{3\_gauss.py}:
\begin{center}
\texttt{7 | 0.3500 | 1.9536351488 | 1.9539537925}
\end{center}

Прямоугольники на этой же строке дали \(1{,}95385677\) при \(N = 8\).
Три ответа расходятся не больше чем на \(0{,}00032\).
Это меньше допуска \(0{,}001\).

Столбца \(N\) у Гаусса нет: сетку не мельчат.
В таблицу лабораторной для этой строки пишутся три числа:
прямоугольники \(1{,}95385677\) с \(N = 8\),
Гаусс по трём узлам \(1{,}9536351488\),
Гаусс по четырём узлам \(1{,}9539537925\).

\end{document}
